\documentclass[a4paper,12pt]{article}

% -------------------------------------------------
% Кодировка и язык
% -------------------------------------------------
\usepackage[T2A]{fontenc}
\usepackage[utf8]{inputenc}
\usepackage[russian]{babel}

% -------------------------------------------------
% Математика
% -------------------------------------------------
\usepackage{amsmath,amssymb,mathtools}
\usepackage{icomma}

% -------------------------------------------------
% Шрифт и типографика
% -------------------------------------------------
\usepackage{helvet}
\renewcommand{\familydefault}{\sfdefault}
\usepackage{microtype}
\usepackage{setspace}
\onehalfspacing
\usepackage{parskip}

% -------------------------------------------------
% Вёрстка
% -------------------------------------------------
\usepackage[
    a4paper,
    left=20mm,
    right=20mm,
    top=20mm,
    bottom=22mm
]{geometry}

\usepackage{enumitem}
\usepackage{tabularx,booktabs}
\usepackage{fancyhdr}
\usepackage{xcolor}
\usepackage{tikz}
\usetikzlibrary{calc}
\usepackage{hyperref}

% -------------------------------------------------
% Палитра
% -------------------------------------------------
\definecolor{accent}{HTML}{008080}
\definecolor{darkaccent}{HTML}{004D4D}
\definecolor{textgray}{HTML}{333333}
\definecolor{softgray}{HTML}{707070}
\definecolor{lightaccent}{HTML}{DCEEEE}

\color{textgray}

% -------------------------------------------------
% Ссылки
% -------------------------------------------------
\hypersetup{
    colorlinks=true,
    linkcolor=darkaccent,
    urlcolor=darkaccent,
    citecolor=darkaccent
}

% -------------------------------------------------
% Колонтитулы
% -------------------------------------------------
\pagestyle{fancy}
\fancyhf{}
\fancyhead[L]{\small\color{softgray}Математика, 6 класс}
\fancyhead[R]{\small\color{softgray}Надя Клименко}
\fancyfoot[C]{\small\color{softgray}\thepage}
\renewcommand{\headrulewidth}{0.3pt}
\renewcommand{\footrulewidth}{0pt}

% -------------------------------------------------
% Заголовки и служебные команды
% -------------------------------------------------
\newcommand{\sectionline}{%
    \vspace{0.2em}
    {\color{accent}\rule{\linewidth}{0.6pt}}
    \vspace{0.5em}
}

\newcommand{\checkitem}{\(\square\)\;}

\newcommand{\important}[1]{%
    \par\medskip
    {\color{darkaccent}\textbf{#1}}
    \par\smallskip
}

\newcommand{\exampletitle}[1]{%
    \par\medskip
    {\color{accent}\textbf{Пример. #1}}
    \par\smallskip
}

\begin{document}

% =================================================
% ТИТУЛЬНЫЙ ЛИСТ
% =================================================

\begin{titlepage}
    \thispagestyle{empty}

    \begin{center}
        \vspace*{25mm}

        {\Large\color{softgray}\textbf{ЧЕК-ЛИСТ ПО МАТЕМАТИКЕ}}

        \vspace{14mm}

        {\fontsize{25}{31}\selectfont
        \color{darkaccent}\textbf{Проценты, средняя скорость\\[2mm]
        и наибольший общий делитель}}

        \vspace{8mm}

        {\large
        Повторение вычислительных приёмов\\
        и переход к теме НОД}

        \vfill

        {\large
        \textbf{Лёвин Артём Александрович}\\[2mm]
        эксклюзивно для \textbf{Нади Клименко}}

        \vspace{9mm}

        {\large \today}

        \vspace{25mm}
    \end{center}

    % Сдержанная кривая Безье
    \begin{tikzpicture}[remember picture,overlay]
        \draw[
            accent,
            line width=1.2pt
        ]
        ($(current page.south west)+(18mm,20mm)$)
        .. controls
        ($(current page.south west)+(60mm,31mm)$)
        and
        ($(current page.south east)+(-65mm,10mm)$)
        ..
        ($(current page.south east)+(-18mm,22mm)$);
    \end{tikzpicture}

\end{titlepage}

% =================================================
% ОСНОВНОЙ ЧЕК-ЛИСТ
% =================================================

\section*{\color{darkaccent}1. Что важно уметь после занятия}
\sectionline

\begin{itemize}[leftmargin=7mm,itemsep=3pt]
    \item \checkitem правильно находить среднюю скорость по всему пути;
    \item \checkitem различать \textbf{целое}, \textbf{часть} и \textbf{процент};
    \item \checkitem находить процент от числа;
    \item \checkitem находить всё число по известной части и её проценту;
    \item \checkitem определять, сколько процентов одна величина составляет от другой;
    \item \checkitem быстро умножать и делить десятичные дроби на \(10\), \(100\), \(1000\);
    \item \checkitem переводить смешанную дробь в неправильную;
    \item \checkitem понимать смысл наибольшего общего делителя;
    \item \checkitem находить НОД разложением чисел на простые множители;
    \item \checkitem находить НОД двух и нескольких чисел.
\end{itemize}

% =================================================
\section*{\color{darkaccent}2. Средняя скорость}
\sectionline

Если движение состоит из нескольких участков, средняя скорость рассчитывается по
\textbf{всему пройденному пути} и \textbf{всему затраченному времени}:

\[
\boxed{
v_{\text{ср}}
=
\frac{S_{\text{общ}}}{t_{\text{общ}}}
}
\]

где

\[
S_{\text{общ}}=S_1+S_2+\ldots,
\qquad
t_{\text{общ}}=t_1+t_2+\ldots
\]

\important{Проверка единиц измерения}

Если скорость требуется получить в километрах в час, расстояние должно быть выражено
в километрах, а время --- в часах.

\[
1\text{ ч}=60\text{ мин},
\qquad
33\text{ мин}
=
\frac{33}{60}\text{ ч}
=
\frac{11}{20}\text{ ч}.
\]

\exampletitle{Движение по трём участкам}

Пусть велосипедист проехал:

\[
1{,}2\text{ км за }6\text{ мин},
\quad
5{,}3\text{ км за }12\text{ мин},
\quad
2{,}3\text{ км за }15\text{ мин}.
\]

Общий путь:

\[
S_{\text{общ}}
=
1{,}2+5{,}3+2{,}3
=
8{,}8\text{ км}.
\]

Общее время:

\[
t_{\text{общ}}
=
6+12+15
=
33\text{ мин}
=
\frac{11}{20}\text{ ч}.
\]

Тогда

\[
v_{\text{ср}}
=
8{,}8:\frac{11}{20}
=
8{,}8\cdot\frac{20}{11}
=
16\text{ км/ч}.
\]

\important{Типичная ошибка}

Нельзя просто находить среднее арифметическое скоростей на отдельных участках.
Надёжная формула:

\[
\boxed{\text{весь путь}\;:\;\text{всё время}}
\]

% =================================================
\section*{\color{darkaccent}3. Проценты: три основных типа задач}
\sectionline

\[
1\%=\frac{1}{100},
\qquad
p\%=\frac{p}{100}.
\]

Главная идея: перед вычислением определите, что в задаче является
\textbf{целым}, а что --- \textbf{частью}.

% -------------------------------------------------
\subsection*{\color{accent}3.1. Найти часть, если известно целое}

Если известно всё число \(A\) и требуется найти \(p\%\) от него:

\[
\boxed{
B=A\cdot\frac{p}{100}
}
\]

\exampletitle{Найти \(25\%\) от \(340\)}

\[
340\cdot\frac{25}{100}
=
85.
\]

Ответ:

\[
\boxed{85}
\]

\medskip


% -------------------------------------------------
\subsection*{\color{accent}3.2. Найти целое, если известны часть и процент}

Если \(B\) составляет \(p\%\) от неизвестного числа \(A\), то

\[
B=A\cdot\frac{p}{100}.
\]

Отсюда

\[
\boxed{
A=B\cdot\frac{100}{p}
}
\]

или

\[
\boxed{
A=B:p\cdot100
}
\]

\exampletitle{\(2\%\) посетителей --- это \(8\) человек}

Известна часть:

\[
8\text{ человек}.
\]

Известен соответствующий процент:

\[
2\%.
\]

Следовательно,

\[
A
=
8\cdot\frac{100}{2}
=
400.
\]

Ответ:

\[
\boxed{400\text{ человек}}
\]

\important{Ключевой ориентир}

Если \textbf{целое неизвестно}, зато известны и величина части, и процент этой части,
используйте

\[
\boxed{\text{часть}\cdot100:\text{процент}}
\]

% -------------------------------------------------
\exampletitle{Задача про пирожные}

Пирожные содержат \(14\%\) сахара. Для их приготовления израсходовали
\(21\) кг сахара.

Сначала найдём общую массу пирожных:

\[
21\cdot\frac{100}{14}
=
150\text{ кг}.
\]

Если масса одного пирожного равна \(100\) г, то

\[
100\text{ г}=0{,}1\text{ кг}.
\]

Количество пирожных:

\[
150:0{,}1=1500.
\]

Ответ:

\[
\boxed{1500\text{ пирожных}}
\]

% -------------------------------------------------
\subsection*{\color{accent}3.3. Найти, сколько процентов составляет часть от целого}

Если известны часть \(B\) и целое \(A\), искомый процент равен

\[
\boxed{
p=\frac{B}{A}\cdot100\%
}
\]

\exampletitle{Яблоки в смеси}

В смеси:

\[
130\text{ г}+270\text{ г}+200\text{ г}=600\text{ г}.
\]

Масса яблок --- \(200\) г.

Следовательно,

\[
p
=
\frac{200}{600}\cdot100\%
=
\frac{100}{3}\%
=
33\frac13\%.
\]

Ответ:

\[
\boxed{33\frac13\%}
\]

% =================================================
\section*{\color{darkaccent}4. Как выбрать формулу в задаче на проценты}
\sectionline

\renewcommand{\arraystretch}{1.45}

\begin{center}
\begin{tabularx}{\linewidth}{@{}>{\bfseries}p{0.28\linewidth}X X@{}}
\toprule
Что известно & Что требуется найти & Формула \\
\midrule

Целое и процент
&
Часть
&
\(\displaystyle
\text{часть}
=
\text{целое}\cdot\frac{p}{100}
\)
\\

Часть и её процент
&
Целое
&
\(\displaystyle
\text{целое}
=
\text{часть}\cdot\frac{100}{p}
\)
\\

Часть и целое
&
Процент
&
\(\displaystyle
p
=
\frac{\text{часть}}{\text{целое}}\cdot100\%
\)
\\

\bottomrule
\end{tabularx}
\end{center}

\important{Полезный вопрос перед решением}

Спросите себя:

\begin{center}
\emph{«Что здесь является всем объектом, а что --- его частью?»}
\end{center}

После этого выбор формулы становится значительно проще.

% =================================================
\section*{\color{darkaccent}5. Процент меньше, результат больше}
\sectionline

Нельзя сравнивать только сами проценты, если они берутся от разных чисел.

Например:

\[
20\%\text{ от }100
=
20,
\]

а

\[
5\%\text{ от }1000
=
50.
\]

Хотя

\[
5\%<20\%,
\]

получается

\[
50>20.
\]

Причина: проценты взяты от разных целых.

\[
\boxed{
\text{результат зависит и от процента, и от исходного числа}
}
\]

% =================================================
\section*{\color{darkaccent}6. Быстрые вычисления с \(10\), \(100\), \(1000\)}
\sectionline

При делении на \(10\), \(100\), \(1000\) десятичная запятая смещается
\textbf{влево} соответственно на \(1\), \(2\), \(3\) разряда:

\[
5{,}6:10=0{,}56,
\]

\[
5{,}6:100=0{,}056,
\]

\[
12{,}34:1000=0{,}01234.
\]

Если цифр слева недостаточно, дописываются нули.

\important{При умножении направление противоположное}

\[
5{,}6\cdot10=56,
\]

\[
5{,}6\cdot100=560,
\]

\[
5{,}6\cdot1000=5600.
\]

Запятая смещается вправо.

% -------------------------------------------------
\subsection*{\color{accent}6.1. Связь с умножением на \(0{,}1\)}

\[
0{,}1=\frac{1}{10},
\]

поэтому

\[
a\cdot0{,}1=a:10.
\]

Например,

\[
560\cdot0{,}1=56.
\]

Аналогично:

\[
a\cdot0{,}01=a:100,
\]

\[
a\cdot0{,}001=a:1000.
\]

% -------------------------------------------------
\subsection*{\color{accent}6.2. Полезный пример с процентами}

\[
10\%\text{ от }5{,}6
=
5{,}6\cdot\frac{10}{100}
=
5{,}6:10
=
0{,}56.
\]

Также

\[
0{,}1\%\text{ от }560
=
560\cdot\frac{0{,}1}{100}
=
560\cdot0{,}001
=
0{,}56.
\]

Результаты совпали:

\[
\boxed{0{,}56}
\]

% =================================================
\section*{\color{darkaccent}7. Смешанные дроби в вычислениях}
\sectionline

Перед умножением и делением смешанную дробь удобно переводить в неправильную.

Общее правило:

\[
a\frac{b}{c}
=
\frac{ac+b}{c}.
\]

Например,

\[
2\frac{1}{21}
=
\frac{2\cdot21+1}{21}
=
\frac{43}{21}.
\]

Если

\[
43\%
\]

некоторого числа равны

\[
2\frac{1}{21},
\]

то

\[
x
=
2\frac{1}{21}\cdot\frac{100}{43}
=
\frac{43}{21}\cdot\frac{100}{43}
=
\frac{100}{21}
=
4\frac{16}{21}.
\]

\important{При делении на дробь}

Деление заменяем умножением на обратную дробь:

\[
a:\frac{b}{c}
=
a\cdot\frac{c}{b}.
\]

После этого обязательно ищем возможность сократить множители.

% =================================================
\newpage

\section*{\color{darkaccent}8. Наибольший общий делитель}
\sectionline

\textbf{Наибольший общий делитель} нескольких натуральных чисел ---
это наибольшее натуральное число, на которое каждое из данных чисел
делится без остатка.

Обозначение:

\[
\operatorname{НОД}(a,b).
\]

Например:

\[
\operatorname{НОД}(12,18)=6.
\]

Действительно, общие делители чисел \(12\) и \(18\):

\[
1,\ 2,\ 3,\ 6,
\]

и наибольший из них --- \(6\).

% -------------------------------------------------
\subsection*{\color{accent}8.1. Метод перебора делителей}

Для небольших чисел можно выписать все делители.

Делители числа \(12\):

\[
1,\ 2,\ 3,\ 4,\ 6,\ 12.
\]

Делители числа \(18\):

\[
1,\ 2,\ 3,\ 6,\ 9,\ 18.
\]

Общие:

\[
1,\ 2,\ 3,\ 6.
\]

Следовательно,

\[
\boxed{\operatorname{НОД}(12,18)=6}
\]

Метод понятный, однако для больших чисел становится громоздким.

% -------------------------------------------------
\subsection*{\color{accent}8.2. Разложение на простые множители}

Более универсальный способ:

\begin{enumerate}[leftmargin=8mm]
    \item разложить каждое число на простые множители;
    \item найти множители, которые присутствуют \textbf{во всех} разложениях;
    \item взять каждый общий простой множитель в наименьшей степени;
    \item перемножить выбранные множители.
\end{enumerate}

\exampletitle{\(\operatorname{НОД}(24,32)\)}

Разложим числа:

\[
24=2^3\cdot3,
\]

\[
32=2^5.
\]

Общий простой множитель --- \(2\).

Минимальная степень двойки:

\[
2^3.
\]

Следовательно,

\[
\boxed{
\operatorname{НОД}(24,32)=2^3=8
}
\]

% -------------------------------------------------
\exampletitle{\(\operatorname{НОД}(18,42)\)}

\[
18=2\cdot3^2,
\]

\[
42=2\cdot3\cdot7.
\]

Общие множители:

\[
2,\quad3.
\]

Поэтому

\[
\operatorname{НОД}(18,42)
=
2\cdot3
=
6.
\]

\[
\boxed{6}
\]

% -------------------------------------------------
\subsection*{\color{accent}8.3. НОД трёх чисел}

Для нескольких чисел действует то же правило:
множитель должен присутствовать \textbf{в каждом} разложении.

\exampletitle{\(\operatorname{НОД}(12,15,18)\)}

\[
12=2^2\cdot3,
\]

\[
15=3\cdot5,
\]

\[
18=2\cdot3^2.
\]

Единственный простой множитель, присутствующий во всех трёх числах, --- \(3\).

Следовательно,

\[
\boxed{
\operatorname{НОД}(12,15,18)=3
}
\]

% =================================================
\section*{\color{darkaccent}9. Признаки делимости, полезные для разложения}
\sectionline

\begin{center}
\begin{tabularx}{\linewidth}{@{}cX@{}}
\toprule
\textbf{Делитель} & \textbf{Признак} \\
\midrule

\(2\)
&
Последняя цифра числа чётная:
\(0,2,4,6,8\).
\\

\(3\)
&
Сумма цифр числа делится на \(3\).
\\

\(5\)
&
Последняя цифра --- \(0\) или \(5\).
\\

\(9\)
&
Сумма цифр числа делится на \(9\).
\\

\(10\)
&
Последняя цифра --- \(0\).
\\

\bottomrule
\end{tabularx}
\end{center}

Эти признаки позволяют быстрее выполнять разложение числа на простые множители.

% =================================================
\section*{\color{darkaccent}10. Типичные ошибки}
\sectionline

\begin{enumerate}[leftmargin=8mm,itemsep=5pt]

    \item
    \textbf{Средняя скорость:}
    использовать среднее арифметическое скоростей вместо формулы
    \[
    \frac{\text{весь путь}}{\text{всё время}}.
    \]

    \item
    \textbf{Единицы времени:}
    подставлять минуты в формулу, если скорость требуется в км/ч.

    \item
    \textbf{Проценты:}
    начинать вычисления до определения целого и части.

    \item
    \textbf{Поиск целого:}
    умножать известную часть на \(\frac{p}{100}\).
    При поиске целого используется обратная операция:
    \[
    \text{часть}\cdot\frac{100}{p}.
    \]

    \item
    \textbf{Поиск процента:}
    менять местами часть и целое.
    Верный порядок:
    \[
    \frac{\text{часть}}{\text{целое}}\cdot100\%.
    \]

    \item
    \textbf{Десятичные дроби:}
    смещать запятую вправо при делении на \(10\), \(100\), \(1000\).

    \item
    \textbf{НОД:}
    включать множитель, который встречается только у части чисел.

    \item
    \textbf{НОД со степенями:}
    брать максимальную степень общего множителя.
    Для НОД выбирается \textbf{минимальная} степень.
\end{enumerate}

% =================================================
\section*{\color{darkaccent}11. Мини-алгоритмы}
\sectionline

\subsection*{\color{accent}Если спрашивают: «Найдите \(p\%\) от числа»}

\[
\boxed{
\text{целое}\cdot\frac{p}{100}
}
\]

\subsection*{\color{accent}Если спрашивают: «Найдите число, если \(p\%\) равны \(\ldots\)»}

\[
\boxed{
\text{известная часть}\cdot\frac{100}{p}
}
\]

\subsection*{\color{accent}Если спрашивают: «Сколько процентов составляет часть?»}

\[
\boxed{
\frac{\text{часть}}{\text{целое}}\cdot100\%
}
\]

\subsection*{\color{accent}Если спрашивают среднюю скорость}

\[
\boxed{
\frac{\text{общий путь}}{\text{общее время}}
}
\]

\subsection*{\color{accent}Если требуется НОД}

\[
\boxed{
\begin{array}{c}
\text{разложить числа на простые множители}\\[1mm]
\downarrow\\[1mm]
\text{оставить множители, общие для всех чисел}\\[1mm]
\downarrow\\[1mm]
\text{взять их в наименьших степенях}\\[1mm]
\downarrow\\[1mm]
\text{перемножить}
\end{array}
}
\]

% =================================================
% ТРЕНИРОВОЧНЫЙ БЛОК
% =================================================

\newpage

\section*{\color{darkaccent}Тренировочный блок}
\sectionline

\textbf{Надя, выполните задания самостоятельно.}
Старайтесь сначала определить тип задачи и только затем выбирать вычисления.

\medskip

\textbf{1. Средняя скорость.}

Велосипедист проехал \(2{,}4\) км за \(9\) минут,
затем \(3{,}6\) км за \(12\) минут и ещё \(4\) км за \(9\) минут.

Найдите среднюю скорость велосипедиста на всём пути в км/ч.

\vspace{8mm}

\textbf{2. Процент от числа.}

Найдите:

\begin{enumerate}[label=\alph*),leftmargin=10mm]
    \item \(18\%\) от \(450\);
    \item \(10\%\) от \(7{,}3\);
    \item \(0{,}1\%\) от \(820\).
\end{enumerate}

\vspace{8mm}

\textbf{3. Нахождение целого и процента.}

\begin{enumerate}[label=\alph*),leftmargin=10mm]
    \item \(12\%\) некоторого числа равны \(36\). Найдите это число.
    \item В коробке \(84\) красных карандаша, всего в коробке \(240\) карандашей.
    Какой процент всех карандашей составляют красные?
\end{enumerate}

\vspace{8mm}

\textbf{4. Проценты в практической задаче.}

Конфеты содержат \(16\%\) орехов.
Для изготовления партии конфет использовали \(24\) кг орехов.

\begin{enumerate}[label=\alph*),leftmargin=10mm]
    \item Найдите общую массу партии конфет.
    \item Сколько конфет изготовили, если масса одной конфеты равна \(50\) г?
\end{enumerate}

\vspace{8mm}

\textbf{5. Наибольший общий делитель.}

Разложите числа на простые множители и найдите:

\begin{enumerate}[label=\alph*),leftmargin=10mm]
    \item \(\operatorname{НОД}(36,60)\);
    \item \(\operatorname{НОД}(48,72)\);
    \item \(\operatorname{НОД}(30,45,75)\).
\end{enumerate}

\vfill

{\small\color{softgray}
Перед проверкой ответов ещё раз убедитесь, что в задачах на проценты правильно
определены целое и часть, а в задаче на среднюю скорость время переведено в часы.
}

% =================================================
% ОТВЕТЫ
% =================================================

\newpage

\section*{\color{darkaccent}Ответы к тренировочному блоку}
\sectionline

\renewcommand{\arraystretch}{1.5}

\begin{center}
\begin{tabularx}{\linewidth}{@{}cX@{}}
\toprule
\textbf{№} & \textbf{Ответ} \\
\midrule

1
&
\(\displaystyle 20\text{ км/ч}\)
\\

2
&
а) \(81\);
\qquad
б) \(0{,}73\);
\qquad
в) \(0{,}82\)
\\

3
&
а) \(300\);
\qquad
б) \(35\%\)
\\

4
&
а) \(150\text{ кг}\);
\qquad
б) \(3000\) конфет
\\

5
&
а) \(12\);
\qquad
б) \(24\);
\qquad
в) \(15\)
\\

\bottomrule
\end{tabularx}
\end{center}

\vspace{12mm}

\section*{\color{darkaccent}Самопроверка}
\sectionline

После выполнения упражнений отметьте:

\begin{itemize}[leftmargin=7mm,itemsep=5pt]
    \item \checkitem Я различаю целое и часть в задачах на проценты.
    \item \checkitem Я знаю три основные формулы для задач на проценты.
    \item \checkitem Я умею быстро делить и умножать на \(10\), \(100\), \(1000\).
    \item \checkitem Я считаю среднюю скорость через весь путь и всё время.
    \item \checkitem Я умею раскладывать натуральное число на простые множители.
    \item \checkitem Я понимаю, какие множители необходимо включать в НОД.
\end{itemize}

\vfill


\end{document}